% R591: EGA I, Chapter0, whole section3.3 heading and3.3.1; source-faithful translation.
% Traducción: OpenAI Codex, GPT-6.1 Sol, esfuerzo Ultra.
\subsection{Pegado de haces}
\label{subsection:0.3.3-fr}

\begin{env}[3.3.1]
\label{0.3.3.1-fr}
Supongamos, como antes, que la categoría $\boldsymbol{K}$ admite límites proyectivos
(generalizados). Sean $X$ un espacio topológico y
$\mathfrak{U}=(U_\lambda)_{\lambda\in L}$ un cubrimiento abierto de $X$; para
cada $\lambda\in L$, sea $\mathcal{F}_\lambda$ un haz sobre
$U_\lambda$, con valores en $\boldsymbol{K}$~; para todo par de índices
$(\lambda,\mu)$, supongamos dado un \emph{isomorfismo}
\[
  \theta_{\lambda\mu}:
  \mathcal{F}_\mu|(U_\lambda\cap U_\mu)
  \xrightarrow{\sim}
  \mathcal{F}_\lambda|(U_\lambda\cap U_\mu);
\]
además, supongamos que, para toda terna $(\lambda,\mu,\nu)$, al denotar por
$\theta'_{\lambda\mu}$, $\theta'_{\mu\nu}$, $\theta'_{\lambda\nu}$ las
restricciones de $\theta_{\lambda\mu}$, $\theta_{\mu\nu}$,
$\theta_{\lambda\nu}$ a $U_\lambda\cap U_\mu\cap U_\nu$, se tenga
\[
  \theta'_{\lambda\nu}
  =\theta'_{\lambda\mu}\circ\theta'_{\mu\nu}
\]
(\emph{condición de pegado} para los $\theta_{\lambda\mu}$). Entonces,
existe un haz $\mathcal{F}$ sobre $X$, con valores en $\boldsymbol{K}$, y para
cada $\lambda$ un isomorfismo
\[
  \eta_\lambda:\mathcal{F}|U_\lambda\xrightarrow{\sim}\mathcal{F}_\lambda
\]
tales que, para todo par $(\lambda,\mu)$, al denotar por
$\eta'_\lambda$ y $\eta'_\mu$ las restricciones de $\eta_\lambda$ y
$\eta_\mu$ a $U_\lambda\cap U_\mu$, se tenga
\[
  \theta_{\lambda\mu}=\eta'_\lambda\circ{\eta'_\mu}^{-1};
\]
además, $\mathcal{F}$ y los $\eta_\lambda$ quedan determinados, por estas
condiciones, hasta un único isomorfismo. La unicidad se deduce, en efecto,
inmediatamente de (3.2.5). Para probar la existencia de $\mathcal{F}$, denotemos por
$\mathfrak{B}$ la base de abiertos formada por los abiertos contenidos en al menos un
$U_\lambda$ y, para cada $U\in\mathfrak{B}$, elijamos (mediante la
función $\boldsymbol{\tau}$ de Hilbert) uno de los $\mathcal{F}_\lambda(U)$ para uno de los
$\lambda$ tales que $U\subset U_\lambda$~; si se denota este objeto por
$\mathcal{F}(U)$, los $\rho_U^V$ para $U\subset V$, $U\in\mathfrak{B}$,
$V\in\mathfrak{B}$ se definen de manera evidente (mediante los
$\theta_{\lambda\mu}$), y la condición de transitividad se sigue de la
condición de pegado~; además, la verificación de
(F\textsubscript{0}) es inmediata; por tanto, el prehaz sobre $\mathfrak{B}$
así definido es efectivamente un haz, y de él se obtiene mediante el procedimiento general
(3.2.1) un haz (ordinario), también denotado por $\mathcal{F}$, que resuelve el
problema. Se dice que $\mathcal{F}$ se obtiene por \emph{pegado de los
$\mathcal{F}_\lambda$ mediante los $\theta_{\lambda\mu}$} y habitualmente se identificarán
$\mathcal{F}_\lambda$ y $\mathcal{F}|U_\lambda$ mediante
$\eta_\lambda$.

Es claro que todo haz $\mathcal{F}$ sobre $X$ con valores en
$\boldsymbol{K}$ puede considerarse obtenido por pegado de los haces
$\mathcal{F}_\lambda=\mathcal{F}|U_\lambda$ (donde $(U_\lambda)$ es un
cubrimiento abierto arbitrario de $X$), mediante los isomorfismos
$\theta_{\lambda\mu}$ que se reducen a la identidad.
\end{env}
